\documentclass[11pt,a4paper]{article}
\usepackage[T1]{fontenc}
\usepackage[utf8]{inputenc}
\usepackage{newtxtext}
\usepackage{amsmath,amsthm,mathtools}
\usepackage{newtxmath}
\usepackage[margin=27mm,headheight=15pt]{geometry}
\usepackage{microtype}
\usepackage{booktabs,tabularx,array}
\usepackage[dvipsnames]{xcolor}
\usepackage{enumitem}
\usepackage{fancyhdr}
\usepackage{hyperref}
\usepackage{bookmark}
\hypersetup{colorlinks=true,linkcolor=MidnightBlue,citecolor=MidnightBlue,
 urlcolor=MidnightBlue,pdftitle={Affine Monodromy and Spherical Residues in Quaternionic Analysis},
 pdfauthor={Research draft prepared with ChatGPT},
 pdfsubject={Global Fueter inversion, integrable spherical singularities, and a sharp energy law}}
\urlstyle{same}
\numberwithin{equation}{section}
\newtheorem{theorem}{Theorem}[section]
\newtheorem{proposition}[theorem]{Proposition}
\newtheorem{lemma}[theorem]{Lemma}
\newtheorem{corollary}[theorem]{Corollary}
\theoremstyle{definition}
\newtheorem{definition}[theorem]{Definition}
\newtheorem{example}[theorem]{Example}
\theoremstyle{remark}
\newtheorem{remark}[theorem]{Remark}
\newcommand{\R}{\mathbb R}
\newcommand{\C}{\mathbb C}
\newcommand{\HH}{\mathbb H}
\newcommand{\HC}{\mathbb H_{\C}}
\newcommand{\SSS}{\mathbb S}
\newcommand{\Z}{\mathbb Z}
\newcommand{\SR}{\mathcal{SR}}
\newcommand{\AM}{\mathcal{AM}}
\newcommand{\D}{\mathcal D}
\newcommand{\lap}{\Delta_4}
\newcommand{\lift}{\mathcal I}
\newcommand{\Per}{\operatorname{Per}}
\newcommand{\Rea}{\operatorname{Re}}
\newcommand{\Ima}{\operatorname{Im}}
\newcommand{\ind}{\operatorname{wind}}
\newcommand{\dd}{\,\mathrm d}
\newcommand{\ds}{\,\mathrm ds}
\newcommand{\supp}{\operatorname{supp}}
\newcommand{\dist}{\operatorname{dist}}
\newcommand{\ii}{\mathbf i}
\newcommand{\jj}{\mathbf j}
\newcommand{\kk}{\mathbf k}
\newcommand{\norm}[1]{\lVert#1\rVert}
\newcommand{\abs}[1]{\lvert#1\rvert}
\newcommand{\bp}{\boldsymbol\beta}
\setlength{\parskip}{2pt}
\setlength{\emergencystretch}{2em}
\setlist{itemsep=3pt,topsep=4pt}
\pagestyle{fancy}
\fancyhf{}
\fancyhead[L]{\small Affine monodromy and spherical residues}
\fancyhead[R]{\small\thepage}
\renewcommand{\headrulewidth}{0.3pt}
\title{\vspace{-10mm}\textbf{Affine Monodromy and Spherical Residues\\in Quaternionic Analysis}\\[5pt]
\large Global Fueter inversion, integrable singularities,\\and a sharp logarithmic energy law}
\author{Research draft prepared with ChatGPT}
\date{21 September 2026}
\begin{document}
\maketitle
\thispagestyle{empty}
\begin{abstract}
The local inverse of the quaternionic Fueter map has an affine ambiguity.
We turn that ambiguity into an explicit global obstruction theory. For an
axial left-monogenic field $g(x+Iy)=P(x,y)+IQ(x,y)$, two simultaneously
closed quaternion-valued one-forms have zero periods precisely when $g$
has a single-valued slice-regular Fueter primitive on the same domain.
On a planar base with $m$ independent holes, reflected logarithmic modes
split the obstruction map and give a cokernel of real dimension $8m$.
We then prove an isolated-sphere singularity theorem: every axial monogenic
field locally integrable across a nonreal two-sphere is, modulo a smooth
monogenic field, a unique combination of two quaternionic residue modes.
The same residues are the affine monodromy coefficients, the coefficients
of a surface-supported Fueter source, and the parameters in an exact
logarithmic $L^2$-energy law. In particular, finite local $L^2$ mass is
removable, whereas nonremovable examples exist in every $L^p$, $1\le p<2$.
A Cauchy--Fueter layer formula gives canonical decaying representatives and
classifies all decaying, locally integrable axial fields with finitely many
spherical singularities. All the stated results are proved; symbolic and
numerical checks are supplied separately.
\end{abstract}
\noindent\textbf{Keywords.} Fueter map; slice regularity; affine monodromy;
period obstruction; spherical residue; removable singularity; surface source.

\medskip
\noindent\textbf{Research status.} The local inverse Fueter theorem is established
literature, not a new result here. The proposed contributions are the explicit
global obstruction--residue correspondence and its classification and energy
consequences. These formulations were not located in the sources checked for
this draft, but originality is provisional: the literature search was not
exhaustive and the manuscript has not undergone independent peer review.
No previously named open conjecture is claimed to have been settled.
\newpage
\begingroup
\small
\setlength{\parskip}{0pt}
\tableofcontents
\endgroup
\newpage

\section{The research problem and the distinction between local and global}
\label{sec:intro}

The Fueter transform sends a slice-regular function $f$ to the left-monogenic
function $\lap f$. Its local inverse is not unique: one can add $qa+b$, with
$a,b\in\HH$, without changing the image. This familiar kernel suggests a
specific question that is not answered merely by constructing a primitive in
small rectangles:
\begin{quote}
When do local Fueter primitives glue to one single-valued slice-regular
function on the original domain, and what analytic information is carried
by their failure to glue?
\end{quote}
The supplied expository manuscript~\cite{Uploaded} ends its inverse construction
with two successive integrations and explicitly leaves nonzero periods on
multiply connected planar sets untreated. We start from that construction,
replace the second integration form by a primitive-independent one, and pursue
the resulting obstruction all the way to a classification of integrable
spherical singularities.

There are three related, but different, issues. Existence in a small
neighborhood is a differential question. Existence on a fixed multiply
connected domain is a single-valuedness question. Removability across a
missing sphere is a regularity question. The same two quaternionic numbers
per hole will govern the second issue and, with a local integrability
hypothesis, the third. Without that hypothesis, zero periods need not imply
removability.

\subsection{What is established background}
The stem formulation of slice regularity is developed by Ghiloni and
Perotti~\cite{GP}. Inverse Fueter integral representations were obtained by
Colombo, Sabadini, and Sommen~\cite{CSS}. Colombo, Pe\~na Pe\~na, Sabadini,
and Sommen~\cite{CPSS} gave a further constructive formula; their
Theorem~3 is explicitly written on a rectangular planar base bounded away
from the real axis. Those results supply the local analytic background for
this paper, not its proposed novelty.

Surjectivity is also discussed in Dong and Qian~\cite{DQ}, including a
contour formulation in their Theorem~3.6. It is important not to interpret
broad surjectivity language as automatically asserting a globally
single-valued primitive on every prescribed domain with holes. Our criterion
specifies that interpretation exactly and supplies explicit counterexamples
to unrestricted same-domain surjectivity. A claim that a particular published
theorem is incorrect would require a separate audit of all its branch and
domain conventions; that claim is not made here.

The August 2026 preprint of De Martino and Pinton~\cite{DMP} concerns inverses
of factorizations of the Fueter map and generalized polyanalytic extensions.
Its questions concerning the intermediate operators are different from our
period problem. We do not present their results, or the questions they
already answer, as new discoveries of this manuscript.

\subsection{Main conclusions}
For $g=P+IQ$, the two forms introduced below are
\begin{equation}\label{eq:intro-forms}
 \alpha_g=\frac{-P\dd x+Q\dd y}{2},\qquad
 \beta_g=\frac{(xP+yQ)\dd x+(yP-xQ)\dd y}{2}.
\end{equation}
Their periods are the coefficients of the affine increment $qa+b$ acquired
by a local primitive around a loop. Their vanishing is necessary and
sufficient for a global primitive. Explicit logarithmic representatives
realize every period vector on a finitely connected base
(Theorems~\ref{thm:period} and~\ref{thm:exact}).

Near the sphere $S_a=u+v\SSS$, $v>0$, an $L^1$ axial monogenic field has a
unique pair $(h,k)\in\HH^2$ and a decomposition
\begin{equation}\label{eq:intro-decomp}
 g=\Psi_a h+\Phi_a k+g_0,
\end{equation}
where $g_0$ extends smoothly through $S_a$. The model fields $\Phi_a,\Psi_a$
are explicit and are regular everywhere else, including the real axis
(Theorem~\ref{thm:L1}). Their residues coincide with the periods in
\eqref{eq:intro-forms}. In the sense of distributions,
\begin{equation}\label{eq:intro-source}
 \D g=\sigma_a\delta_{S_a},\qquad
 \sigma_a(p)=\frac{2}{v}(ph+k),\quad p\in S_a.
\end{equation}
Here $\delta_{S_a}$ means surface measure, not a point mass.

Finally, for fixed sufficiently small $R$, as $\varepsilon\downarrow0$,
\begin{equation}\label{eq:intro-energy}
 \int_{\varepsilon<\dist(q,S_a)<R}|g(q)|^2\dd V_q
 =8\bigl(|uh+k|^2+v^2|h|^2\bigr)
       \log\frac{R}{\varepsilon}+E_R+o(1).
\end{equation}
The finite constant $E_R$ depends on the field, not just its residues.
The logarithmic coefficient is also
$(2\pi)^{-1}\norm{\sigma_a}_{L^2(S_a)}^2$.
This connects a topological obstruction, a local singularity invariant,
and a quantitative analytic divergence.

\section{Conventions and the axial differential calculus}
\label{sec:prelim}

\subsection{Domains, stems, and multiplication order}
Write
\[
 q=x_0+\ii x_1+\jj x_2+\kk x_3,
 \qquad \ii^2=\jj^2=\kk^2=\ii\jj\kk=-1.
\]
Quaternionic coefficients will always be on the right. The sphere of
imaginary units is $\SSS=\{I\in\HH:I^2=-1\}$. A nonreal quaternion is
uniquely $q=x+Iy$ with $y=|\Ima q|>0$.

Let $U\subset\C_+=\{x+\iota y:y>0\}$ be a connected open set. The symbol
$\iota$ denotes a central complex imaginary unit, not one of $\ii,\jj,\kk$.
Set
\[
 \Omega_U=\{x+Iy:x+\iota y\in U,\ I\in\SSS\}.
\]
This is a product domain, disjoint from the real axis. A holomorphic map
$F=A+\iota B:U\to\HC$, where $A,B:U\to\HH$, induces
\[
 \lift(F)(x+Iy)=A(x,y)+IB(x,y).
\]
Equivalently, extend the stem to $U\cup\bar U$ by
$F(\bar z)=\kappa(F(z))$, where $\kappa(A+\iota B)=A-\iota B$.
We denote the resulting space of left slice-regular functions by $\SR(U)$.
Holomorphy means, componentwise,
\begin{equation}\label{eq:CR}
 A_x=B_y,\qquad A_y=-B_x.
\end{equation}
The definition on product domains is the stem definition; no uniqueness
principle requiring a real interval is being imported into it.

The left Fueter operator and its conjugate are
\[
 \D=\partial_{x_0}+\ii\partial_{x_1}+\jj\partial_{x_2}+\kk\partial_{x_3},
 \qquad
 \overline\D=\partial_{x_0}-\ii\partial_{x_1}-\jj\partial_{x_2}-\kk\partial_{x_3}.
\]
Their product in either order is $\lap$. We write $\AM(U)$ for the smooth
axial fields $g=P+IQ$ on $\Omega_U$ that satisfy $\D g=0$. The functions
$P,Q$ may be quaternion valued; left monogenicity does not assert right
monogenicity for such general coefficients.

\subsection{The differential identities}
\begin{lemma}[Axial calculus]\label{lem:axial}
For smooth quaternion-valued $A,B$ and $y>0$,
\begin{align}
 \D(A+IB)&=\left(A_x-B_y-\frac{2B}{y}\right)
                 +I(B_x+A_y),\label{eq:axial-D}\\
 \lap(A+IB)&=A_{xx}+A_{yy}+\frac{2A_y}{y}
    +I\left(B_{xx}+B_{yy}+\frac{2B_y}{y}-\frac{2B}{y^2}\right).
 \label{eq:axial-lap}
\end{align}
Consequently $g=P+IQ$ belongs to $\AM(U)$ precisely when
\begin{equation}\label{eq:Vekua}
 P_x-Q_y-\frac{2Q}{y}=0,\qquad Q_x+P_y=0.
\end{equation}
For $f=A+IB\in\SR(U)$,
\begin{equation}\label{eq:Fueter}
 \lap f=\frac{2A_y}{y}
       +I\left(\frac{2B_y}{y}-\frac{2B}{y^2}\right),
 \qquad \D\lap f=0.
\end{equation}
\end{lemma}
\begin{proof}
Let $w=\Ima q$ and $I=w/y$. With $e_1=\ii,e_2=\jj,e_3=\kk$,
\[
 \partial_{x_j}y=x_j/y,\qquad
 \partial_{x_j}I=e_j/y-wx_j/y^3.
\]
Thus $\sum e_j\partial_{x_j}I=-2/y$, the three-dimensional Laplacian of
$I$ is $-2I/y^2$, and
$\sum (\partial_{x_j}I)(\partial_{x_j}y)=0$. The product rule gives
\eqref{eq:axial-D} and \eqref{eq:axial-lap}. Evaluating at $I$ and $-I$
shows that the two quaternionic coefficients in \eqref{eq:axial-D} vanish
independently, which gives \eqref{eq:Vekua}.

The Cauchy--Riemann equations make $A$ and $B$ planar harmonic, proving the
first formula in \eqref{eq:Fueter}. If its two coefficients are $P,Q$, then
\[
 P_x=\frac{2B_{yy}}y=Q_y+\frac{2Q}{y},\qquad
 Q_x+P_y=\frac{2(B_{xy}+A_{yy})}{y}
          -\frac{2(B_x+A_y)}{y^2}=0.
\]
This proves the last assertion and is the quaternionic axial Fueter
mapping theorem in the conventions used here.
\end{proof}

\begin{lemma}[Affine kernel]\label{lem:kernel}
On a connected $U\subset\C_+$,
\[
 \ker(\lap:\SR(U)\longrightarrow\AM(U))
       =\{q\mapsto qh+k:h,k\in\HH\}.
\]
\end{lemma}
\begin{proof}
If \eqref{eq:Fueter} is zero, then $A_y=0$ and $yB_y=B$.
Set $C=B/y$. Then $C_y=0$ and, by \eqref{eq:CR},
$C_x=B_x/y=-A_y/y=0$. Therefore $C=h$ is constant. Now
$A_x=B_y=h$ and $A_y=0$, so $A=xh+k$. Connectedness makes both constants
global. The converse follows by applying the Laplacian to $qh+k$.
\end{proof}

\subsection{An axial norm and the four-dimensional measure}
For an axial field define
\[
 N_g(x,y)=\bigl(|P(x,y)|^2+|Q(x,y)|^2\bigr)^{1/2}.
\]
If $d\omega$ is the ordinary area measure on the unit two-sphere, then
\begin{equation}\label{eq:sphere-average}
 \frac1{4\pi}\int_{\SSS}|P+IQ|^2\dd\omega(I)=|P|^2+|Q|^2.
\end{equation}
Indeed, the cross term in the squared norm is linear in the three real
coordinates of $I$ and has zero average; $|IQ|=|Q|$. In particular,
\begin{equation}\label{eq:norm-equiv}
 N_g\le\sup_I|g(x+Iy)|\le\sqrt2\,N_g.
\end{equation}
Polar coordinates only in the imaginary three-space give
\begin{equation}\label{eq:measure}
 \dd V=y^2\dd x\dd y\dd\omega(I).
\end{equation}
These identities retain the noncommutative field values while reducing
certain norm questions to the planar base.

\section{Two simultaneous period obstructions}
\label{sec:periods}

\begin{lemma}[The two closed forms]\label{lem:closed}
For every $g=P+IQ\in\AM(U)$, the forms $\alpha_g,\beta_g$ in
\eqref{eq:intro-forms} are closed. If $g=\lap f$ for $f=A+IB\in\SR(U)$,
then
\begin{equation}\label{eq:potentials}
 \alpha_g=\dd C,\qquad \beta_g=\dd D_0,
 \qquad C=\frac By,\quad D_0=A-\frac{xB}{y}.
\end{equation}
\end{lemma}
\begin{proof}
The coefficient of $\dd x\wedge\dd y$ in $\dd\alpha_g$ is
$(Q_x+P_y)/2$. The corresponding coefficient for $\beta_g$ is
\[
 \frac12\{y(P_x-Q_y)-2Q-x(Q_x+P_y)\}.
\]
Both vanish by \eqref{eq:Vekua}.

For a Fueter image, \eqref{eq:Fueter} and \eqref{eq:CR} imply
\[
 C_x=-P/2,\qquad C_y=Q/2.
\]
Differentiating $D_0=A-xC$ gives
\[
 (D_0)_x=(xP+yQ)/2,\qquad (D_0)_y=(yP-xQ)/2,
\]
which proves \eqref{eq:potentials}.
\end{proof}

\begin{theorem}[Global single-valued inversion]\label{thm:period}
Let $U\subset\C_+$ be connected and open, and let $g\in\AM(U)$.
The following conditions are equivalent:
\begin{enumerate}[label=\textup{(\roman*)}]
\item There is a single-valued $f\in\SR(U)$ with $\lap f=g$.
\item Both forms $\alpha_g,\beta_g$ are exact on $U$.
\item For every closed piecewise $C^1$ curve $\gamma$ in $U$,
\[
 \int_\gamma\alpha_g=\int_\gamma\beta_g=0.
\]
\end{enumerate}
When these conditions hold, choose primitives $C,D_0$ of the two forms.
Then the formula
\begin{equation}\label{eq:global-inverse}
 \boxed{\ f(x+Iy)=(x+Iy)C(x,y)+D_0(x,y)\ }
\end{equation}
is a Fueter primitive. Every other primitive is $f+qh+k$ for unique
$h,k\in\HH$.
\end{theorem}
\begin{proof}
The implication (i)$\Rightarrow$(ii) is Lemma~\ref{lem:closed}.
Exact forms have zero periods. Conversely, a closed form with zero
integrals on all closed curves has a global primitive obtained by path
integration from a base point. Path independence follows by joining one
path to the reverse of another. This argument applies separately to the
four real components of each quaternion-valued form, proving
(iii)$\Rightarrow$(ii).

For (ii)$\Rightarrow$(i), let $A=xC+D_0$ and $B=yC$. The displayed
formulas for $\dd C,\dd D_0$ yield
\begin{align*}
 A_x&=C+x(-P/2)+(xP+yQ)/2=C+yQ/2=B_y,\\
 A_y&=xQ/2+(yP-xQ)/2=yP/2=-B_x.
\end{align*}
Thus $A+\iota B$ is holomorphic. Moreover,
\[
 \frac{2A_y}{y}=P,\qquad
 \frac{2B_y}{y}-\frac{2B}{y^2}
       =\frac{2(C+yQ/2)}y-\frac{2yC}{y^2}=Q.
\]
Equation~\eqref{eq:Fueter} proves $\lap f=g$. The ambiguity is exactly
Lemma~\ref{lem:kernel}.
\end{proof}

The second form is not the second form in a successive integration scheme.
If $C$ is first chosen with $\dd C=\alpha_g$, that scheme uses
\[
 \eta_C=(C+yQ/2)\dd x+(yP/2)\dd y.
\]
A direct identity is
\begin{equation}\label{eq:remove-primitive}
 \eta_C-\dd(xC)=\beta_g.
\end{equation}
Thus $\beta_g$ eliminates the dependence on a first primitive. Both
obstructions can be measured directly from $g$, before solving either
integration problem.

\begin{corollary}[Affine monodromy]\label{cor:monodromy}
A local Fueter primitive can be continued along every path in $U$.
After continuation around a based loop $\gamma$, its increment is
\begin{equation}\label{eq:monodromy}
 q\,a_\gamma+b_\gamma,\qquad
 a_\gamma=\int_\gamma\alpha_g,\quad
 b_\gamma=\int_\gamma\beta_g.
\end{equation}
The map $\gamma\mapsto(a_\gamma,b_\gamma)$ is an additive homomorphism
and factors through $H_1(U;\Z)$.
\end{corollary}
\begin{proof}
Integrate both closed forms along a lifted path and use
\eqref{eq:global-inverse}. Local changes in the two primitives are constant
quaternions; their effect on $f$ is $qa+b$. For a closed path the constants
are exactly the two integrals. Additivity is additivity of path integrals.
Closed-form integrals vanish on boundaries by Stokes' theorem, so they
factor through first homology.
\end{proof}

\begin{corollary}[A useful normalization]\label{cor:normalization}
Fix $z_*=x_*+\iota y_*\in U$. Whenever $g$ has a global primitive,
there is exactly one primitive that vanishes on the entire sphere
$x_*+y_*\SSS$. It is obtained by taking $C(z_*)=D_0(z_*)=0$.
\end{corollary}
\begin{proof}
The construction gives existence. An affine function $qh+k$ vanishing
for all $q=x_*+Iy_*$ has, by subtracting the values at $I$ and $-I$,
$2Iy_*h=0$. Thus $h=0$ and then $k=0$.
\end{proof}

\subsection{Extending the criterion through the real axis}
The period criterion is not limited to examples disjoint from the real
axis. The following extension will be used for punctured quaternionic
space.

\begin{proposition}[Real-axis extension]\label{prop:axis}
Let $D\subset\C$ be open and conjugation invariant, and suppose
$U=D\cap\C_+$ is connected. Let $g$ be a smooth axial left-monogenic
field on the full circularization $\Omega_D$, including any real points.
Then $g$ has a single-valued slice-regular primitive on $\Omega_D$ if and
only if its two periods vanish on $U$. The affine kernel is unchanged.
\end{proposition}
\begin{proof}
Necessity follows by restriction. For sufficiency, construct $C,D_0$ on
$U$ by Theorem~\ref{thm:period}. Near a real point $x_0\in D$, fix
$J\in\SSS$ and define, also for negative $y$,
\[
 P(x,y)=\frac{g(x+Jy)+g(x-Jy)}2,\qquad
 Q(x,y)=-\frac J2\bigl(g(x+Jy)-g(x-Jy)\bigr).
\]
These agree with the original coefficients for $y>0$ and are smooth,
with $P$ even and $Q$ odd. The forms $\alpha_g,\beta_g$ extend smoothly
across $y=0$. They remain closed there by continuity. Their pullbacks
under $(x,y)\mapsto(x,-y)$ equal themselves.

On a small symmetric disk about $x_0$, take primitives of these extended
forms. Each primitive is even: its difference from its reflection is
constant, and that constant is zero at a real point. Add constant
quaternions so that these primitives agree with $C,D_0$ on the upper
half-disk. Then $A=xC+D_0$ is even, $B=yC$ is odd, and the same calculation
as in Theorem~\ref{thm:period} proves the Cauchy--Riemann equations across
the diameter. These local stems therefore give a slice-regular extension.
On overlaps they agree with the original upper stem and consequently
agree everywhere. Reflecting the upper stem handles the remaining lower
part of $D$. The equality $\lap f=g$ extends by smoothness. Uniqueness
modulo $qh+k$ follows on $U$ and then throughout $\Omega_D$.
\end{proof}

\section{Explicit logarithmic representatives of the obstruction}
\label{sec:modes}

\subsection{Why the logarithms are reflected}
Fix $a=u+\iota v\in\C_+$ and its sphere $S_a=u+v\SSS$. Locally away
from $a,\bar a$, put
\begin{equation}\label{eq:Lambda}
 \Lambda_a(z)=\frac{\log(z-a)-\log(z-\bar a)}{2\pi\iota}.
\end{equation}
On the upper half-plane the second logarithm can be chosen globally,
since $z-\bar a$ has positive imaginary part. Continuation around $a$
changes $\Lambda_a$ by $1$ and $z\Lambda_a$ by $z$.
Define locally
\begin{equation}\label{eq:mode-definition}
 \Phi_a=\lap\lift(\Lambda_a),\qquad
 \Psi_a=\lap\lift(z\Lambda_a).
\end{equation}
Changes of logarithmic branch are real constants in $\Lambda_a$ and
real affine terms in $z\Lambda_a$. The affine kernel makes the two
images single valued.

The reflected logarithm is not needed for that single-valuedness on
$\C_+$. It has two additional advantages: its images extend through the
real axis, and they decay at infinity. It also makes the logarithm in
the explicit formulas a dimensionless ratio.

\begin{proposition}[Canonical modes and their periods]\label{prop:modes}
The fields $\Phi_a,\Psi_a$ are single-valued smooth axial monogenic
functions on $\HH\setminus S_a$. For every closed curve
$\gamma\subset\C_+\setminus\{a\}$,
\begin{align}
 \int_\gamma\alpha_{\Phi_a}&=0,&
 \int_\gamma\beta_{\Phi_a}&=\ind(\gamma,a),\label{eq:period-Phi}\\
 \int_\gamma\alpha_{\Psi_a}&=\ind(\gamma,a),&
 \int_\gamma\beta_{\Psi_a}&=0.\label{eq:period-Psi}
\end{align}
Here winding number uses the counterclockwise-positive convention.
\end{proposition}
\begin{proof}
The preceding branch calculation and Lemma~\ref{lem:kernel} give
single-valuedness off the real axis. Local Fueter mapping gives
monogenicity. Near any real $x_0$, the ratio
$(z-a)/(z-\bar a)$ is nonzero and has modulus $1$ at $x_0$. Choose its
local logarithm purely imaginary at that point. The resulting
$\Lambda_a$ satisfies
$\Lambda_a(\bar z)=\overline{\Lambda_a(z)}$ on a small symmetric disk:
both sides have the same derivative and value at $x_0$.
It therefore induces a slice-regular function across the real axis, as
does $z\Lambda_a$. Their Fueter images supply compatible smooth
extensions.

Continuation around $\gamma$ changes the two local slice primitives
by $\ind(\gamma,a)$ and $q\ind(\gamma,a)$, respectively. The lower
point $\bar a$ has zero winding for a curve in $\C_+$. Comparing with
\eqref{eq:monodromy} proves all four identities.
\end{proof}

\subsection{Branch-free formulas}
Set, for $y>0$,
\begin{align}
 s&=x-u,& d_-&=s^2+(y-v)^2,& d_+&=s^2+(y+v)^2,\label{eq:dpm}\\
 \ell&=\frac12\log\frac{d_-}{d_+},&
 H&=\frac1{d_-}-\frac1{d_+},&
 T&=\frac{y-v}{d_-}-\frac{y+v}{d_+}.\label{eq:HT}
\end{align}
Then
\begin{align}
 \Phi_a(x+Iy)&=P_0+IQ_0,&
 P_0&=\frac{sH}{\pi y},&
 Q_0&=\frac{\ell}{\pi y^2}-\frac{T}{\pi y},\label{eq:Phi-explicit}\\
 \Psi_a(x+Iy)&=P_1+IQ_1,&
 P_1&=\frac{xsH+yT+\ell}{\pi y},&
 Q_1&=\frac{ysH-xT}{\pi y}+\frac{x\ell}{\pi y^2}.
 \label{eq:Psi-explicit}
\end{align}
All coefficients in these formulas are real. Values on the real axis
are their smooth extensions; the apparent powers of $1/y$ should not
be evaluated separately at $y=0$.

\begin{proof}[Derivation of the formulas]
For a single logarithm $L(z)=\log(z-a)/(2\pi\iota)$, write
$\rho=|z-a|$, $t=y-v$, and choose a local argument $\theta$.
Its stem coefficients are
\[
 A=\theta/(2\pi),\qquad B=-\log\rho/(2\pi).
\]
The derivatives $A_y=s/(2\pi\rho^2)$ and
$B_y=-t/(2\pi\rho^2)$ give
\[
 P^\circ_0=\frac{s}{\pi y\rho^2},\qquad
 Q^\circ_0=\frac{\log\rho}{\pi y^2}-\frac{t}{\pi y\rho^2}.
\]
The stem $zL$ has coefficients $xA-yB$ and $yA+xB$. In its Fueter
image the undifferentiated argument cancels, leaving
\[
 P^\circ_1=\frac{xs+yt}{\pi y\rho^2}+\frac{\log\rho}{\pi y},\qquad
 Q^\circ_1=\frac{ys-xt}{\pi y\rho^2}+\frac{x\log\rho}{\pi y^2}.
\]
Subtract the same expressions with $a$ replaced by $\bar a$.
The differences of $\log\rho$, $1/\rho^2$, and $t/\rho^2$ are precisely
$\ell,H,T$, proving \eqref{eq:Phi-explicit}--\eqref{eq:Psi-explicit}.
\end{proof}

\begin{proposition}[Decay]\label{prop:decay}
Let
\begin{equation}\label{eq:E}
 E(q)=\frac{\bar q}{2\pi^2|q|^4},\qquad q\ne0.
\end{equation}
As $|q|\to\infty$,
\begin{equation}\label{eq:mode-decay}
 \Phi_a(q)=8\pi v E(q)+O(|q|^{-4}),\qquad
 \Psi_a(q)=8\pi uv E(q)+O(|q|^{-4}).
\end{equation}
The estimates are uniform in direction.
\end{proposition}
\begin{proof}
On $|z|>|a|$, choose the branch of \eqref{eq:Lambda} tending to zero
at infinity. Its Laurent series is
\begin{equation}\label{eq:Lambda-series}
 \Lambda_a(z)=\sum_{n\ge1}\lambda_n z^{-n},\qquad
 \lambda_n=-\frac{\Ima(a^n)}{\pi n}\in\R,
 \quad \lambda_1=-\frac v\pi,\quad\lambda_2=-\frac{uv}\pi.
\end{equation}
Here $\Ima(a^n)$ means the real coefficient of $\iota$ in the complex
number $a^n$. The induced quaternionic Laurent series and its derivatives
converge uniformly outside every larger radius. Direct differentiation
in four dimensions gives
\[
 \lap(q^{-1})=-\frac{4\bar q}{|q|^4}=-8\pi^2E(q).
\]
For completeness, $q^{-1}=\bar q|q|^{-2}$, the function $|q|^{-2}$ is
harmonic away from zero in dimension four, and the cross term is
$-4\bar q|q|^{-4}$. Taking the leading term in the two Laurent series
proves \eqref{eq:mode-decay}. The constant term of $z\Lambda_a$ has
zero Laplacian. Uniform derivative convergence gives the stated remainders.
\end{proof}

\begin{example}[A genuine failure of unrestricted global inversion]
\label{ex:no-global}
On $\HH\setminus\SSS$, take $a=\iota$ and $g=\Phi_{\iota}$.
This is a globally defined smooth axial monogenic field, even across the
real axis. On a small counterclockwise circle around $\iota$ in the
upper base, its $\beta$-period is $1$. Therefore no single-valued
slice-regular $f$ on $\HH\setminus\SSS$ satisfies $\lap f=g$.
The same example on a small circular annulus about $\iota$ gives a
product-domain counterexample. This is a topological failure, not a
failure of the local inverse theorem.
\end{example}

\section{Finite topology, a split exact sequence, and quantitative inversion}
\label{sec:finite}

\subsection{The topological hypothesis}
Assume that there are points $a_1,\ldots,a_m\in\C_+\setminus U$ and
closed piecewise $C^1$ curves $\gamma_1,\ldots,\gamma_m\subset U$ such that
\begin{equation}\label{eq:winding-basis}
 H_1(U;\Z)\longrightarrow\Z^m,\qquad
 [\gamma]\longmapsto(\ind(\gamma,a_1),\ldots,\ind(\gamma,a_m))
\end{equation}
is an isomorphism and $\ind(\gamma_j,a_k)=\delta_{jk}$.
This explicit hypothesis includes a bounded planar domain with $m$
smooth Jordan holes and closure in $\C_+$: put one point in each hole
and take positively oriented hole loops. Cutting along arcs from the
holes to the outer boundary gives the familiar winding basis. It also
includes $\C_+\setminus\{a_1,\ldots,a_m\}$ and disks with finitely many
punctures. Stating \eqref{eq:winding-basis} avoids imposing unnecessary
boundary regularity on the theorem.

Set
\begin{equation}\label{eq:period-vector}
 a_j(g)=\int_{\gamma_j}\alpha_g,\qquad
 b_j(g)=\int_{\gamma_j}\beta_g,\qquad
 \Per(g)=(a_1(g),b_1(g),\ldots,a_m(g),b_m(g)).
\end{equation}
The functions $a_j(g)$ should not be confused with the fixed missing
points $a_j$; the argument $(g)$ distinguishes the period coefficients.

\begin{theorem}[Explicit finite-dimensional cokernel]\label{thm:exact}
Under \eqref{eq:winding-basis}, the following is an exact sequence of
right quaternionic vector spaces:
\begin{equation}\label{eq:exact-sequence}
 0\longrightarrow\HH^2
 \xrightarrow{(h,k)\mapsto qh+k}\SR(U)
 \xrightarrow{\lap}\AM(U)
 \xrightarrow{\Per}\HH^{2m}\longrightarrow0.
\end{equation}
The period map has the explicit right inverse
\begin{equation}\label{eq:splitting}
 \mathcal S((h_j,k_j)_{j=1}^m)
       =\sum_{j=1}^m(\Psi_{a_j}h_j+\Phi_{a_j}k_j).
\end{equation}
Consequently
\begin{equation}\label{eq:direct-sum}
 \AM(U)=\lap\SR(U)\ \oplus\
   \operatorname{span}_{\HH}\{\Psi_{a_j},\Phi_{a_j}:1\le j\le m\},
\end{equation}
and the cokernel of $\lap$ has quaternionic dimension $2m$, or real
dimension $8m$. In particular, $\lap$ is surjective in this class exactly
when $m=0$.
\end{theorem}
\begin{proof}
Exactness at $\SR(U)$ is Lemma~\ref{lem:kernel}. Closed-form periods
depend only on homology, so vanishing on the chosen winding basis is
vanishing on every closed curve. Theorem~\ref{thm:period} therefore gives
$\ker\Per=\lap\SR(U)$.

Equations~\eqref{eq:period-Phi} and \eqref{eq:period-Psi} show
$\Per\mathcal S=\operatorname{id}_{\HH^{2m}}$. This proves surjectivity
of the final map and linear independence of the displayed modes modulo
the image. For an arbitrary $g$, the field
$g-\mathcal S\Per(g)$ has zero periods and is a Fueter image. A vector
in the intersection of the two summands has both zero periods and its
own coefficient vector as its periods, so it is zero. All assertions follow.
\end{proof}

If the field is originally defined on a circular domain meeting the real
axis, and its upper base satisfies \eqref{eq:winding-basis}, the same
exact sequence holds for the full-domain spaces. Indeed, the correction
modes already extend across the real axis, and
Proposition~\ref{prop:axis} extends the primitive of the zero-period
remainder. In particular, the theorem applies on
$\HH\setminus\bigcup_{j=1}^m S_{a_j}$.

\subsection{Continuous projection and path bounds}
Define
\begin{equation}\label{eq:projection}
 \Pi g=\mathcal S\Per(g).
\end{equation}
Then $\Pi^2=\Pi$, and $\ker\Pi=\lap\SR(U)$. This is a concrete
projection, not just an existence statement about a complement.

\begin{proposition}[Stability of the obstruction and reconstruction]
\label{prop:stability}
Let $\gamma$ have length $L_\gamma$ and
$R_\gamma=\max_{z\in\gamma}|z|$. Then
\begin{equation}\label{eq:period-bound}
 \left|\int_\gamma\alpha_g\right|
       \le\frac{L_\gamma}{2}\sup_\gamma N_g,\qquad
 \left|\int_\gamma\beta_g\right|
       \le\frac{L_\gamma R_\gamma}{2}\sup_\gamma N_g.
\end{equation}
The image $\lap\SR(U)$ is closed in the compact-open topology of
$\AM(U)$, and also in its usual $C^\infty$ topology.

For zero-period $g$, normalize the primitive at $z_*$. Suppose every point
of a compact $K\subset U$ can be reached from $z_*$ by a path of length
at most $L$ contained in a compact $K'\subset U$, and put
$R=\max_{K'}|z|$. Then
\begin{equation}\label{eq:inverse-bound}
 \sup_{q\in\Omega_K}|f(q)|\le LR\sup_{K'}N_g.
\end{equation}
Such a path family exists for every compact $K$, with suitable $K',L$.
The normalized inverse on the zero-period subspace is continuous in
$C^\infty$ on compact subsets.
\end{proposition}
\begin{proof}
For a unit planar tangent $(\dot x,\dot y)$, the norm of
$-P\dot x+Q\dot y$ is at most $N_g$. The real coefficient pair in
$2\beta_g(\dot\gamma)$ is
$(x\dot x+y\dot y,y\dot x-x\dot y)$; its Euclidean norm is $|z|$.
Cauchy--Schwarz in $\HH^2\cong\R^8$ gives \eqref{eq:period-bound}.
Thus each period is continuous on compact sets. The range is the
intersection of their kernels, even without finite topology, by
Theorem~\ref{thm:period}; it is closed.

Along any one of the prescribed paths,
\[
 |C(z)|\le\tfrac L2\sup_{K'}N_g,\qquad
 |D_0(z)|\le\tfrac {LR}2\sup_{K'}N_g.
\]
Since $|q|=|z|\le R$, formula~\eqref{eq:global-inverse} proves
\eqref{eq:inverse-bound}. To find a uniform path family, cover $K$ by
finitely many small disks compactly contained in $U$, connect their
centers to $z_*$ by finitely many polygonal paths, and append short
segments in the disks. Their compact union and slightly larger disks
give $K'$, and their lengths have a common finite bound.

All first derivatives of $C,D_0$ are the explicit coefficients in
\eqref{eq:intro-forms}. Higher derivatives are derivatives of those
coefficients. Together with the path estimate, these observations give
compact-set $C^k$ estimates for $f$ in terms of finitely many $C^{k-1}$
seminorms of $g$ (with the zeroth-order bound handled separately).
This proves the stated $C^\infty$ continuity.
\end{proof}

\begin{corollary}[A certificate of non-approximability]\label{cor:distance}
If a compact $K\subset U$ contains all the curves $\gamma_j$, then
\begin{equation}\label{eq:distance}
 \inf_{h\in\lap\SR(U)}\sup_K N_{g-h}
 \ge \max_{1\le j\le m}
 \left\{\frac{2|a_j(g)|}{L_j},
        \frac{2|b_j(g)|}{L_jR_j}\right\},
\end{equation}
where $L_j$ and $R_j$ are the lengths and maximal radii of the curves.
An explicit approximant from the range is $g-\Pi g$, so
\begin{equation}\label{eq:distance-upper}
 \inf_{h\in\lap\SR(U)}\sup_K N_{g-h}
 \le \sum_j\left(
  \sup_K N_{\Psi_{a_j}}\,|a_j(g)|+
  \sup_K N_{\Phi_{a_j}}\,|b_j(g)|\right).
\end{equation}
\end{corollary}
\begin{proof}
Every $h$ in the range has zero periods. Apply \eqref{eq:period-bound}
to $g-h$ for the lower bound. For the upper bound use
$g-\Pi g\in\lap\SR(U)$ and the triangle inequality in $\HH^2$.
\end{proof}

\section{Integrable singularities along a two-sphere}
\label{sec:singularities}

Fix $a=u+\iota v$, with $v>0$, and choose $0<R<v$. Put
\[
 B_a(R)=\{z:|z-a|<R\},\qquad
 \mathcal T_a(R)=\Omega_{B_a(R)}.
\]
For $q=x+Iy\in\mathcal T_a(R)$, define
\begin{equation}\label{eq:rho}
 \rho=\sqrt{(x-u)^2+(y-v)^2}=\dist(q,S_a),\qquad
 p(q)=u+Iv\in S_a.
\end{equation}
The point $p(q)$ is the nearest point of the sphere to $q$. A punctured
planar disk therefore describes a genuine tubular puncture in four real
dimensions.

\subsection{Integrability and the leading term}
For fixed $P,Q$, the identities
\[
 P=\frac1{4\pi}\int_{\SSS}(P+IQ)\dd\omega,
 \qquad Q=-\frac1{4\pi}\int_{\SSS}I(P+IQ)\dd\omega
\]
imply that $g$ is $L^1$ across $S_a$ if and only if $P,Q$ are $L^1$
across $a$ in the planar base. For one direction use these averages and
$y\ge v-R>0$ in \eqref{eq:measure}; for the other use
$|P+IQ|\le|P|+|Q|$ and $y\le v+R$. The analogous equivalence holds for
any finite $p\ge1$, by Jensen's inequality.

\begin{lemma}[Principal part of the residue modes]\label{lem:principal}
Let $h,k\in\HH$, $c=uh+k$, and
$G=\Psi_a h+\Phi_a k=P+IQ$. Uniformly in the normal angle and in
$I\in\SSS$, as $\rho\downarrow0$,
\begin{align}
 P&=\frac{cs+vh\,t}{\pi v\rho^2}+O(1+|\log\rho|),\label{eq:P-principal}\\
 Q&=\frac{vh\,s-ct}{\pi v\rho^2}+O(1+|\log\rho|),\label{eq:Q-principal}
\end{align}
where $s=x-u$, $t=y-v$. Equivalently,
\begin{equation}\label{eq:geometric-principal}
 G(q)=\frac1{\pi v}(q-p(q))^{-1}(p(q)h+k)
                        +O(1+|\log\rho|).
\end{equation}
In particular, $G$ belongs to $L^p(\mathcal T_a(R))$ for every
$1\le p<2$.
\end{lemma}
\begin{proof}
In the single-log formulas used to derive
\eqref{eq:Phi-explicit}--\eqref{eq:Psi-explicit}, replace the smooth
coefficients $x,y,1/y$ in terms of order $\rho^{-1}$ by $u,v,1/v$.
Each error is $O(1)$ because the coefficient changes by $O(\rho)$.
The remaining logarithmic terms are $O(1+|\log\rho|)$.
The reflected logarithm has no singularity at $a$ and contributes a
smooth term. This gives the two displayed coefficients, with right
quaternionic multiplication retained.

Since $q-p(q)=s+It$, its inverse is $(s-It)/\rho^2$. Multiplication by
$c+Ivh$ gives exactly the numerators of
\eqref{eq:P-principal} and \eqref{eq:Q-principal}. Finally the normal
area element is $\rho\dd\rho\dd\theta$, while $y$ is bounded above
and below. The bound $O(\rho^{-1}+|\log\rho|)$ is $L^p$ exactly in the
asserted subcritical range (with no need to assert nonintegrability yet).
\end{proof}

\subsection{A distributional defect with only two coefficients}
Define the real-linear planar operator
\begin{equation}\label{eq:L}
 \mathscr L(P,Q)=\left(P_x-Q_y-\frac{2Q}{y},\ Q_x+P_y\right).
\end{equation}
It is a first-order operator with smooth coefficients on $B_a(R)$.
We regard pairs as elements of $\HH^2$ and interpret distributions
componentwise over $\R$.

\begin{lemma}[An $L^1$ puncture creates only a point mass]
\label{lem:defect}
Suppose $P,Q\in L^1(B_a(R))$ and $\mathscr L(P,Q)=0$ away from $a$.
Then there are unique $s_0,s_1\in\HH$ such that, distributionally,
\begin{equation}\label{eq:defect}
 \mathscr L(P,Q)=(s_0,s_1)\delta_a.
\end{equation}
No derivatives of $\delta_a$ can occur.
\end{lemma}
\begin{proof}
The defect is supported at the single point $a$. A distribution
supported at a point is a finite sum of derivatives of the point mass;
a proof is included in Appendix~\ref{app:point}. If the highest nonzero
derivative order were $r\ge1$, choose a test function $\varphi$ whose
$r$-jet detects a nonzero top-order coefficient and put
\[
 \varphi_\varepsilon(z)=\varepsilon^r
                    \varphi((z-a)/\varepsilon).
\]
The pairing of the defect with this test tends to a nonzero constant;
terms of lower order tend to zero. On the other hand, the first-order
form of $\mathscr L$ and boundedness of $1/y$ on the disk give
\[
 |\langle\mathscr L(P,Q),\varphi_\varepsilon\rangle|
 \le C\bigl(\varepsilon^{r-1}+\varepsilon^r\bigr)
       \int_{|z-a|<C'\varepsilon}(|P|+|Q|)\dd x\dd y.
\]
The right side tends to zero by absolute integrability, including when
$r=1$. This is a contradiction. Only the point-mass term remains.
\end{proof}

\begin{lemma}[Defects of the canonical modes]\label{lem:mode-defect}
For $G=\Psi_a h+\Phi_a k$ and $c=uh+k$,
\begin{equation}\label{eq:mode-defect}
 \mathscr L(P_G,Q_G)=\left(\frac{2c}{v},\ 2h\right)\delta_a.
\end{equation}
\end{lemma}
\begin{proof}
Introduce the central complexification $W=P_G+\iota Q_G$. The complex
version of $\mathscr L$ is
\[
 2\bar\partial W-2Q_G/y,\qquad
 \bar\partial=\tfrac12(\partial_x+\iota\partial_y).
\]
Lemma~\ref{lem:principal} gives
\begin{equation}\label{eq:W-principal}
 W=\frac{c+\iota vh}{\pi v(z-a)}+O(1+|\log|z-a||).
\end{equation}
There is only a point-mass defect by Lemma~\ref{lem:defect}. Its
coefficient is computed by integration around a small circle. The
boundary term for $2\bar\partial$ is
$\int(n_x+\iota n_y)W\ds$. With
$z-a=\varepsilon e^{\iota\theta}$, its principal part integrates to
$2(c+\iota vh)/v$. The boundary contribution of the logarithmic
remainder is $O(\varepsilon(1+|\log\varepsilon|))$ and vanishes.
The integral of the lower-order term $2Q_G/y$ over a shrinking disk
also vanishes because it is $L^1$. Equivalently, this is the elementary
identity $\bar\partial(1/(z-a))=\pi\delta_a$, with its normalization
verified by the same circle computation. Separating the central real
and imaginary coefficients proves the claim.
\end{proof}

\begin{lemma}[Weak axial monogenicity is removable]\label{lem:weak}
If $P,Q\in L^1(B_a(R))$ and $\mathscr L(P,Q)=0$ distributionally on
the full disk, then $P+IQ$ has a smooth axial monogenic representative
on the full tube $\mathcal T_a(R)$.
\end{lemma}
\begin{proof}
The axial differential formula~\eqref{eq:axial-D} remains valid for
$L^1$ coefficients in the sense of distributions on a tube with $y>0$.
One way to see this is to mollify $P,Q$ in the planar variables and use
\eqref{eq:axial-D} for the smooth approximants. Their axial lifts converge
locally in $L^1$ by \eqref{eq:measure}. Pairing the right side with a
four-dimensional test function amounts to pairing the planar coefficients
with its spherical averages multiplied by $y^2$; those are smooth test
functions. Thus the identity passes to the limit.

It follows that the lifted field satisfies $\D g=0$ distributionally.
Applying $\overline\D$ gives $\lap g=0$. The $L^1$ version of Weyl's
lemma, proved in Appendix~\ref{app:Weyl}, makes every real component
smooth and harmonic. Since $\D g=0$ still holds distributionally, it
holds pointwise for that representative. The spherical averages recover
smooth $P,Q$ and agree with the original coefficients almost everywhere.
They preserve axiality through the missing sphere.
\end{proof}

\subsection{The classification theorem}
\begin{theorem}[Complete $L^1$ normal form at an isolated sphere]
\label{thm:L1}
Let $g$ be smooth, axial, and left monogenic on
$\mathcal T_a(R)\setminus S_a$. Assume it is locally integrable across
$S_a$; equivalently, after possibly reducing $R$,
\[
 \int_{\mathcal T_a(R)}|g(q)|\dd V_q<\infty.
\]
Then there are unique $h,k\in\HH$ and a unique smooth axial monogenic
field $g_0$ on the full tube such that
\begin{equation}\label{eq:L1-normal}
 \boxed{\ g=\Psi_a h+\Phi_a k+g_0.\ }
\end{equation}
For every sufficiently small counterclockwise circle $\gamma$ around
$a$ in the planar base,
\begin{equation}\label{eq:residue-period}
 h=\int_\gamma\alpha_g,\qquad k=\int_\gamma\beta_g.
\end{equation}
Thus the space of $L^1$ axial monogenic germs at $S_a$, modulo removable
germs, is a right quaternionic vector space of dimension $2$.
\end{theorem}
\begin{proof}
By the norm equivalence preceding Lemma~\ref{lem:principal}, the planar
coefficients are $L^1$. Lemma~\ref{lem:defect} supplies a defect
$(s_0,s_1)\delta_a$. Set
\begin{equation}\label{eq:source-to-residue}
 h=\frac{s_1}{2},\qquad
 k=\frac{vs_0-us_1}{2}.
\end{equation}
Then $2(uh+k)/v=s_0$. Lemma~\ref{lem:mode-defect} says that
$g-\Psi_a h-\Phi_a k$ has zero planar defect. It is still $L^1$, so
Lemma~\ref{lem:weak} gives the smooth monogenic extension $g_0$.

The forms of a smooth monogenic field on a disk are closed and have
zero circle periods. The four periods of the two modes are given by
\eqref{eq:period-Phi}--\eqref{eq:period-Psi}. Applying them to the
normal form proves \eqref{eq:residue-period}. They determine $h,k$
uniquely, after which $g_0$ is also unique. Conversely every $(h,k)$ is
realized by the corresponding mode combination, which is $L^1$ by
Lemma~\ref{lem:principal}. This proves the dimension assertion.
\end{proof}

\begin{definition}[Affine spherical residue]
The ordered pair
\[
 \operatorname{Res}^{\mathrm{aff}}_{S_a}(g)=(h,k)
\]
from Theorem~\ref{thm:L1} is the affine spherical residue in this
manuscript. The term is explicitly defined here; it is not being
identified with the usual point-singularity Fueter flux residue.
\end{definition}

The contour formulas make sense even without the $L^1$ assumption.
What the theorem adds is completeness: in the $L^1$ class no further
singular invariants survive modulo smooth fields. Section~\ref{sec:examples}
will show why this fails beyond that class.

\subsection{The exact surface source}
For a continuous quaternion-valued density $\sigma$ on $S_a$, define
$\sigma\delta_{S_a}$ by
\[
 \langle\sigma\delta_{S_a},\varphi\rangle
     =\int_{S_a}\sigma(p)\varphi(p)\dd S_p
\]
for real scalar test functions; all quaternionic components are treated
in the usual distributional sense.

\begin{theorem}[Residue equals affine surface source]\label{thm:source}
Under the assumptions of Theorem~\ref{thm:L1}, the $L^1$ field on the
filled tube satisfies
\begin{equation}\label{eq:surface-source}
 \boxed{\ \D g=\sigma_a\delta_{S_a},\qquad
       \sigma_a(p)=\frac2v(ph+k).\ }
\end{equation}
Hence the two period obstructions describe precisely the possible
surface-supported source densities in this axial integrable class.
\end{theorem}
\begin{proof}
The planar defect is
$\bigl(2(uh+k)/v,2h\bigr)\delta_{(u,v)}$ by the normal form and
Lemma~\ref{lem:mode-defect}. Lift the distributional identity as in
Lemma~\ref{lem:weak}. In the volume element
$y^2\dd x\dd y\dd\omega$, the two planar point masses give
\[
 v^2\int_{\SSS}\left(\frac{2(uh+k)}v+2Ih\right)
                    \varphi(u+Iv)\dd\omega(I).
\]
Since $\dd S_p=v^2\dd\omega(I)$ on $S_a$, this is exactly
$\langle(2/v)(ph+k)\delta_{S_a},\varphi\rangle$. The removable field
has no source.
\end{proof}

The source density is an affine function of the point $p$ on the
sphere, with coefficients on the right. It is not an arbitrary function
on that sphere. This finite-dimensional restriction is where axiality
has substantial content; general monogenic fields are not classified
by this theorem.

\section{The residue--energy law and the sharp integrability threshold}
\label{sec:energy}

For the residue pair $(h,k)$, define the positive-definite real quadratic
form
\begin{equation}\label{eq:quadratic}
 \mathcal Q_a(h,k)=|uh+k|^2+v^2|h|^2.
\end{equation}
Because $v>0$, it vanishes exactly when $h=k=0$.

\begin{theorem}[Exact logarithmic coefficient and finite renormalization]
\label{thm:energy}
Under the assumptions of Theorem~\ref{thm:L1}, choose a smaller fixed
$R$ whose closed tube lies in the region of the normal form. Then a
finite real number $E_R(g)$ exists such that
\begin{equation}\label{eq:energy-law}
 \boxed{\
 \int_{\varepsilon<\dist(q,S_a)<R}|g(q)|^2\dd V_q
 =8\mathcal Q_a(h,k)\log\frac R\varepsilon+E_R(g)+o(1).
 \ }
\end{equation}
Equivalently, the logarithmic coefficient is
\begin{equation}\label{eq:source-energy}
 8\mathcal Q_a(h,k)
       =\frac1{2\pi}\int_{S_a}|\sigma_a(p)|^2\dd S_p.
\end{equation}
The coefficient is unchanged by adding a smooth axial monogenic field.
The finite renormalized constant need not be unchanged.
\end{theorem}
\begin{proof}
Let $c=uh+k$. By Theorem~\ref{thm:L1} and
Lemma~\ref{lem:principal}, write $g=G_{-1}+R_0$, where
\[
 G_{-1}(q)=\frac{s-It}{\pi v\rho^2}(c+Ivh),\qquad
 |R_0(q)|\le C(1+|\log\rho|).
\]
Multiplicativity of the quaternionic norm gives
\[
 |G_{-1}(q)|^2=\frac{|c+Ivh|^2}{\pi^2v^2\rho^2}.
\]
The average of $|c+Ivh|^2$ over $\SSS$ is
$|c|^2+v^2|h|^2=\mathcal Q_a(h,k)$, by
\eqref{eq:sphere-average}.

Use $s=\rho\cos\theta$, $t=\rho\sin\theta$ and
$y=v+\rho\sin\theta$. If $y^2$ in the volume element is first replaced
by $v^2$, integration of the leading term gives
\[
 \frac{v^2}{\pi^2v^2}(4\pi)(2\pi)\mathcal Q_a(h,k)
             \int_\varepsilon^R\frac{\dd\rho}{\rho}
 =8\mathcal Q_a(h,k)\log(R/\varepsilon).
\]
The change $y^2-v^2=O(\rho)$ has an absolutely integrable effect at
$\rho=0$. So do all remaining terms: the cross term is bounded in
absolute value by
$C\rho^{-1}(1+|\log\rho|)$, and the square of the remainder by
$C(1+|\log\rho|)^2$. Multiplication by the normal area factor $\rho$
makes both integrable. Therefore the difference between the full
integral and the displayed logarithm converges to a finite limit,
not merely a bounded remainder. This proves \eqref{eq:energy-law}.

Finally $\sigma_a(u+Iv)=2(c+Ivh)/v$ and
$\dd S=v^2\dd\omega$, so
\[
 \int_{S_a}|\sigma_a|^2\dd S
 =4\int_{\SSS}|c+Ivh|^2\dd\omega
 =16\pi\mathcal Q_a(h,k).
\]
This gives \eqref{eq:source-energy} and the final assertions.
\end{proof}

Here ``energy'' means the squared $L^2$ norm of the field, not its
Dirichlet integral $\int|\nabla g|^2$. That distinction matters for the
critical exponent and the logarithm.

\begin{corollary}[Sharp $L^2$ removability]\label{cor:L2}
An axial monogenic field with an isolated spherical singularity that is
locally $L^2$ across the sphere extends smoothly and monogenically
through it. More precisely, in the $L^1$ class the following are equivalent:
\[
 \text{removability}\quad\Longleftrightarrow\quad
 (h,k)=(0,0)\quad\Longleftrightarrow\quad
 \text{finite local }L^2\text{ mass}.
\]
For every $1\le p<2$, there are nonremovable examples locally in $L^p$.
Every nonzero combination $\Psi_a h+\Phi_a k$ is such an example.
\end{corollary}
\begin{proof}
A locally $L^2$ field is $L^1$ on a finite-volume tube by
Cauchy--Schwarz. Theorem~\ref{thm:L1} applies. By
Theorem~\ref{thm:energy}, a nonzero residue pair forces logarithmically
divergent $L^2$ mass. Zero residues give the smooth extension directly.
A smooth extension has finite $L^2$ mass on smaller compact tubes.
The subcritical examples are $L^p$ by Lemma~\ref{lem:principal} and
nonremovable because their residue pairs are nonzero.
\end{proof}

\begin{corollary}[A pointwise threshold]\label{cor:pointwise}
If an axial monogenic field near $S_a$ satisfies
\[
 \sup_{\dist(q,S_a)=\rho}|g(q)|=o(\rho^{-1}),\qquad\rho\downarrow0,
\]
then the singularity is removable. Allowing all $O(\rho^{-1})$ fields
would make this false.
\end{corollary}
\begin{proof}
The hypothesis implies local $L^1$ integrability. If a residue pair
were nonzero, the leading term in Lemma~\ref{lem:principal} would have
nonzero sphere-averaged squared norm after multiplication by $\rho^2$,
whereas the remainder tends to zero after multiplication by $\rho$.
This contradicts the assumed uniform little-oh bound. The canonical
modes furnish the $O(\rho^{-1})$ counterexamples.
\end{proof}

\subsection{Why the ordinary Fueter flux loses information}
\begin{proposition}[Flux and its invisible residue component]
\label{prop:flux}
Give $\partial\mathcal T_a(r)$ the outward quaternionic unit normal $n$.
For every sufficiently small $r>0$,
\begin{equation}\label{eq:flux}
 \int_{\partial\mathcal T_a(r)}n(q)g(q)\dd S_q
       =8\pi v(uh+k).
\end{equation}
Thus the ordinary quaternionic flux detects $uh+k$, but not the full
pair $(h,k)$. In particular, $h\ne0$ and $k=-uh$ give a nonremovable
$L^1$ singularity with zero flux.
\end{proposition}
\begin{proof}
The divergence theorem applied componentwise to $g$ gives
$\int_{\partial V}ng=\int_V\D g$ for smooth fields. Applying it to
an annular tube shows that the displayed flux is independent of $r$.
On its inner limiting tube write
$q=u+r\cos\theta+I(v+r\sin\theta)$. Then
\[
 n=\cos\theta+I\sin\theta,\qquad
 \dd S=(v+r\sin\theta)^2r\dd\theta\dd\omega.
\]
Multiplying the principal part by $n$ on the left cancels its normal
inverse and leaves $(c+Ivh)/(\pi vr)$. The logarithmic remainder has
flux tending to zero. The limiting integral is
\[
 \frac v\pi\int_0^{2\pi}\int_{\SSS}(c+Ivh)
                    \dd\omega\dd\theta=8\pi vc,
\]
where $c=uh+k$. Independence of $r$ proves the exact identity.
The final claim follows from Theorem~\ref{thm:L1} or
Corollary~\ref{cor:L2}.
\end{proof}

\subsection{Coordinate invariance of the coefficient}
A real translation $q'=q-t$ changes the affine increment
$qh+k$ to $q'h+(k+th)$ and the center $u$ to $u-t$.
Hence
\[
 |(u-t)h+(k+th)|^2+v^2|h|^2=\mathcal Q_a(h,k).
\]
The energy coefficient is therefore independent of the choice of real
origin. Formula~\eqref{eq:source-energy} gives its intrinsic geometric
meaning without choosing period coordinates at all.

\section{Cauchy layers and the global classification with finitely many spheres}
\label{sec:layers}

\subsection{A normalized fundamental solution}
The kernel $E$ in \eqref{eq:E} is locally integrable in four dimensions
and satisfies
\begin{equation}\label{eq:fundamental}
 \D E=\delta_0
\end{equation}
in the sense of distributions. To verify the normalization directly,
$\D\bar q=4$ and differentiating $|q|^{-4}$ cancels that contribution
away from zero. On $|q|=r$, with $n=q/r$,
\[
 nE=\frac1{2\pi^2r^3}.
\]
The area of that three-sphere is $2\pi^2r^3$, so its total flux is $1$.
Integration by parts on a punctured ball and passage to the limit
prove \eqref{eq:fundamental}; the error from replacing a test function
by its value at zero is $O(r)$. This also fixes every constant in the
following representation.

\begin{theorem}[Canonical sphere-layer representation]\label{thm:layer}
For $q\notin S_a$ and $h,k\in\HH$,
\begin{equation}\label{eq:layer}
 \boxed{\
 \Psi_a(q)h+\Phi_a(q)k
   =\frac2v\int_{S_a}E(q-p)(ph+k)\dd S_p.\ }
\end{equation}
The field on the right is the unique locally integrable field on
$\HH$, tending to zero at infinity, whose distributional Fueter
source is $(2/v)(ph+k)\delta_{S_a}$. Uniqueness here holds even without
assuming axiality of a competing field.
\end{theorem}
\begin{proof}
Define the right side to be $K(q)$. The integral is smooth and
monogenic away from $S_a$ by differentiation under the integral.
It is locally integrable in all four variables: on a compact set,
Fubini's theorem bounds the integral of $|K|$ by a constant times
\[
 \int_{S_a}\int_{\text{compact set}}|q-p|^{-3}\dd V_q\dd S_p,
\]
which is finite, uniformly in $p$. Formula~\eqref{eq:fundamental}
therefore gives
$\D K=(2/v)(ph+k)\delta_{S_a}$ distributionally. The compactness of
$S_a$ and the bound for $E$ imply $K(q)=O(|q|^{-3})$ at infinity.

By Proposition~\ref{prop:modes}, the left side $G$ is smooth on
$\HH\setminus S_a$, and it is locally $L^1$ across $S_a$ by
Lemma~\ref{lem:principal}. Theorem~\ref{thm:source} gives the same
source as $K$; Proposition~\ref{prop:decay} gives its decay.
Thus $G-K$ is globally $L^1_{\mathrm{loc}}$, distributionally
monogenic, and tends to zero at infinity. Applying $\overline\D$
and Weyl's lemma makes it a smooth entire harmonic field. It is
bounded and therefore constant; the constant is zero. The last step
follows, for example, by the harmonic mean-value formula on arbitrarily
large balls and its derivative estimate, which forces all first
derivatives of a bounded entire harmonic function to vanish.

The same argument applies to the difference of any two locally
integrable fields with the specified source and decay, proving uniqueness
without an axiality hypothesis.
\end{proof}

Although axiality was not required for that uniqueness step, the density
$(2/v)(ph+k)$ is a special affine density. Arbitrary surface densities
would give a much larger class of solutions, outside this paper's
classification.

\subsection{A finite-sphere realization theorem}
Let $a_j=u_j+\iota v_j\in\C_+$ be distinct, and write
$\Sigma=\bigcup_{j=1}^m S_{a_j}$. Distinct upper-base points give disjoint
spheres. Denote by $\mathcal M_{1,0}(\Sigma)$ the right quaternionic
vector space of smooth axial monogenic fields on $\HH\setminus\Sigma$
which are locally integrable across every sphere and satisfy
$|g(q)|\to0$ as $|q|\to\infty$. This is local integrability at finite
points, not an assertion of global $L^1(\HH)$ integrability.

\begin{theorem}[Global classification of decaying integrable fields]
\label{thm:global-class}
The residue map is an isomorphism
\begin{equation}\label{eq:global-iso}
 \mathcal M_{1,0}(\Sigma)\ \cong\ \HH^{2m}.
\end{equation}
Its inverse is the explicit formula
\begin{equation}\label{eq:global-field}
 g(q)=\sum_{j=1}^m\bigl(\Psi_{a_j}(q)h_j+\Phi_{a_j}(q)k_j\bigr)
 =\sum_{j=1}^m\frac2{v_j}\int_{S_{a_j}}
              E(q-p)(ph_j+k_j)\dd S_p.
\end{equation}
In particular, this space has real dimension $8m$. Every such field
actually decays at least as $O(|q|^{-3})$, and more precisely
\begin{equation}\label{eq:far-field}
 g(q)=8\pi E(q)\sum_{j=1}^m v_j(u_jh_j+k_j)+O(|q|^{-4}).
\end{equation}
\end{theorem}
\begin{proof}
Apply Theorem~\ref{thm:L1} in disjoint small tubes to extract unique
residues $(h_j,k_j)$. Subtract the finite sum of canonical modes.
Near each sphere all the other modes are smooth, and the singular part
at that sphere has been removed. The remainder therefore extends to
an entire smooth monogenic field. It tends to zero at infinity, since
both $g$ and the modes do. Harmonic Liouville, as used in
Theorem~\ref{thm:layer}, makes the remainder zero. Conversely every
finite mode sum has the required regularity, local integrability, and
decay, and its residue vector is the specified one. This proves the
isomorphism. The layer expression follows from Theorem~\ref{thm:layer},
and the asymptotic formula follows from \eqref{eq:mode-decay}.
\end{proof}

\begin{corollary}[Global realization without a decay condition]
\label{cor:entire-remainder}
Every smooth axial monogenic field on $\HH\setminus\Sigma$ that is
locally integrable across the spheres has a unique decomposition
\[
 g=g_{\mathrm{ent}}+\sum_j(\Psi_{a_j}h_j+\Phi_{a_j}k_j),
\]
where $g_{\mathrm{ent}}$ is entire axial monogenic. It has a global
slice-regular Fueter primitive on $\HH\setminus\Sigma$ exactly when
all the residue pairs vanish. In that case $g=g_{\mathrm{ent}}$ and
it has a primitive on all of $\HH$.
\end{corollary}
\begin{proof}
The same subtraction argument gives the entire remainder, and periods
give uniqueness. A global primitive forces zero periods. Conversely,
if all pairs vanish, the field is entire. Its upper base is the simply
connected half-plane, so Theorem~\ref{thm:period} gives a primitive
there and Proposition~\ref{prop:axis} extends it to all of $\HH$.
\end{proof}

\begin{corollary}[Additivity of the singular energy]
For disjoint sufficiently small tubes about the $S_{a_j}$,
\[
 \sum_j\int_{\varepsilon<\dist(q,S_{a_j})<R_j}|g(q)|^2\dd V_q
 =8\sum_j\mathcal Q_{a_j}(h_j,k_j)\log(1/\varepsilon)+C+o(1).
\]
\end{corollary}
\begin{proof}
Near each sphere the other modes, and any entire remainder, are smooth.
Apply Theorem~\ref{thm:energy} in each tube and add. The fixed
$\log R_j$ terms are absorbed into $C$.
\end{proof}

The exact sequence of Section~\ref{sec:finite} describes the obstruction
to a primitive. Theorem~\ref{thm:global-class} describes actual decaying
fields, not merely a quotient. These two finite-dimensional objects are
linked by the same explicit representatives, but they answer different
questions.

\section{Examples, necessary hypotheses, and reproducible checks}
\label{sec:examples}

\subsection{One sphere and two independent obstructions}
On $\HH\setminus\SSS$, the field $\Phi_{\iota}$ has residue $(0,1)$,
nonzero flux $8\pi$, and energy coefficient $8$. Its local primitive
acquires the constant $1$ around the puncture. The field $\Psi_{\iota}$
has residue $(1,0)$ and the same energy coefficient $8$, but its flux
is zero. Its local primitive acquires $q$ rather than a constant.
Both fields are locally $L^p$ for $1\le p<2$ and both decay at infinity.
Neither has a single-valued global Fueter primitive on the punctured
space. This exhibits the necessity of both forms in
\eqref{eq:intro-forms}, even in the simplest topology.

\subsection{Zero periods alone do not imply removability}
Take a small punctured disk about $a=u+\iota v\in\C_+$ and the
single-valued holomorphic stem
\[
 F(z)=\frac1{z-a}.
\]
Let $f=\lift(F)$ and $g=\lap f$. Then $g$ is a Fueter image, so
both its periods vanish. Nevertheless it has a nonremovable singularity
along $S_a$. To see this without a general singularity theorem, put
$s=x-u$, $t=y-v$, $\rho^2=s^2+t^2$. The stem coefficients are
$A=s/\rho^2$, $B=-t/\rho^2$, so
\[
 P=-\frac{4st}{y\rho^4},\qquad
 Q=\frac{2(t^2-s^2)}{y\rho^4}+\frac{2t}{y^2\rho^2}.
\]
Their leading squared norm is
\[
 \frac{16s^2t^2+4(t^2-s^2)^2}{v^2\rho^8}
       =\frac4{v^2\rho^4}.
\]
Thus $N_g$ is bounded below by a positive multiple of $\rho^{-2}$
near the puncture. Its planar integral diverges like
$\int_0 \rho^{-1}\dd\rho$. The field is not locally $L^1$, exactly
where Theorem~\ref{thm:L1} does not apply. The distinction between
``no monodromy'' and ``no singularity'' is therefore essential.

\subsection{Removing several independent obstructions}
For a base with two holes and missing points $a_1,a_2$, consider
\[
 g=\lap f_0+\Psi_{a_1}\ii+\Phi_{a_1}\jj
              +\Psi_{a_2}(1+\kk)-2\Phi_{a_2}.
\]
Its period vector is $(\ii,\jj,1+\kk,-2)$ in the corresponding winding
basis. The projection $\Pi$ removes exactly the four displayed mode
terms, irrespective of the unknown primitive $f_0$. This example also
shows that the construction works with genuinely noncommuting right
coefficients; no common complex slice for the coefficients is assumed.

\subsection{Numerical energy example}
The accompanying script takes
\[
 u=\frac25,\quad v=\frac{13}{10},\quad
 h=1+2\ii-\jj+\frac12\kk,\quad
 k=\frac14-\ii+3\jj+2\kk.
\]
For $g=\Psi_a h+\Phi_a k$, direct arithmetic gives
$8\mathcal Q_a(h,k)=181$. With $R=0.3$, the computed integrals are:
\begin{center}
\renewcommand{\arraystretch}{1.17}
\begin{tabular}{@{}rrrr@{}}
\toprule
$\varepsilon$ & $E(\varepsilon,R)$ &
$E(\varepsilon,R)-181\log(R/\varepsilon)$ & $J(\varepsilon)$\\
\midrule
$10^{-2}$ & $626.9031484$ & $11.28642233$ & $181.2151049$\\
$10^{-3}$ & $1043.8028211$ & $11.41819314$ & $181.0049681$\\
$10^{-4}$ & $1460.5735515$ & $11.42102173$ & $181.0000892$\\
$10^{-5}$ & $1877.3415023$ & $11.42107075$ & $181.0000014$\\
\bottomrule
\end{tabular}
\end{center}
Here $E(\varepsilon,R)$ is the integral in
\eqref{eq:energy-law}, and the shell density is
\[
 J(\rho)=4\pi\int_0^{2\pi}
       y^2\bigl(|P|^2+|Q|^2\bigr)\rho^2\dd\theta,
 \quad x=u+\rho\cos\theta,\ y=v+\rho\sin\theta.
\]
Thus $E(\varepsilon,R)=\int_{\log\varepsilon}^{\log R}J(e^t)\dd t$.
The convergence of $J$ to $181$ and of the subtracted integral to a
finite limit checks the normalizations numerically. The table uses
160-point Gauss--Legendre rules in both the angular and logarithmic
radial variables; the displayed digits are diagnostic numerical values,
not certified error bounds.

\subsection{What the code does and does not verify}
The file \texttt{verify.py} performs thirteen exact symbolic identities:
the two Vekua equations and closure of both forms for the unreflected
logarithmic modes, their derivation from the Fueter formula, and the
principal quadratic energy identity. Linearity then gives the reflected
modes. It also integrates all sixteen entries of the two-hole period
matrix at 45-decimal working precision; the observed errors in this
run were below $10^{-46}$. A quaternion multiplication check verifies
the spherical norm average on test data. Independent quadrature of the
four-dimensional sphere-layer integral at a real and a nonreal test point
checks the source normalization; its observed errors were below $10^{-12}$.
Finally the script produces the energy
table above and writes the results to \texttt{verification\_results.json}.

These computations are not a formal proof checker. In particular,
they do not prove the distributional defect lemma, the global gluing
theorem, or the exhaustive $L^1$ classification. Those conclusions rest
on the written arguments. Numerical quadrature errors reported by
comparison with the exact periods are observed errors, not interval
certificates. The code is included to make algebraic sign and
normalization checks repeatable, rather than to replace analysis.

\section{Novelty assessment, scope, and remaining research directions}
\label{sec:status}

\subsection{A precise separation of claims}
The differential identities, the affine kernel, local inverse Fueter
existence, and the Cauchy--Fueter fundamental kernel are established
background. No priority is claimed for them. Closed-form exactness,
point-supported distribution structure, and Weyl regularity are also
standard tools; elementary proofs of the needed versions have been
included for completeness.

The candidate-original contribution is the explicit chain of
identifications
\[
 \begin{gathered}
 \text{two primitive-independent period forms}\\
 \Updownarrow\\
 \text{affine monodromy of a Fueter primitive}\\
 \Updownarrow\quad\text{under the local }L^1\text{ hypothesis}\\
 \text{complete two-quaternion spherical residue}\\
 \Updownarrow\\
 \text{affine surface source and its logarithmic }L^2\text{ coefficient}.
 \end{gathered}
\]
Its concrete consequences are the split exact sequence, the explicit
$8m$-dimensional obstruction space, the complete integrable normal form,
and the classification of decaying finite-sphere fields. The formulas
for the source and the positive-definite coefficient specify information
that an ordinary total flux does not capture. Some individual
consequences, especially an $L^2$ removability threshold for a first-order
elliptic equation, may have broader precedents and are not by themselves
the basis of a novelty claim.

The sources actually used in this audit include the primary stem paper,
the inverse Fueter literature cited above, the full 2013 arXiv text of
the 2014 integral-formula paper, the 2018 Dong--Qian preprint, and the
2026 De Martino--Pinton preprint. Bibliographic leads also included
Common--Sommen's axial construction~\cite{Common} and the 2016 inversion
paper~\cite{DKQS}; their complete texts were not available for a full
comparison in this session. The absence of a matching formula from a
limited search is not a proof of absence from published literature.
A specialist literature review remains necessary before claiming
publication priority. This limitation concerns originality assessment,
not an omitted step in any stated theorem.

\subsection{What has and has not been resolved}
The same-domain inverse question is resolved here for the precise
spaces and domains stated: all closed curves must have both periods
zero, and finitely generated winding topology has exactly the explicit
cokernel described above. The integrable isolated-sphere question is
also resolved in the axial class, including the exact source and
energy coefficient. The examples refute the unrestricted assertion
that every axial monogenic field on an arbitrary circular domain must
have a single-valued primitive on that same domain.

No analogous classification is asserted for nonaxial fields. There,
the possible source density on a sphere need not be affine and an
infinite-dimensional analysis is to be expected. No finite-dimensional
cokernel assertion is made for a base with infinitely many independent
holes, although the zero-period criterion remains valid. At a real
point $v=0$, the geometry changes from a two-dimensional singular
sphere to a point; the formulas containing $1/v$ and the threshold
$p=2$ do not apply. Finally, outside the $L^1$ class additional
principal parts can have zero affine residues, as the explicit example
in Section~\ref{sec:examples} demonstrates.

Several extensions are concrete rather than merely terminological.
One may classify higher distributional derivatives of spherical
sources under weaker integrability or prescribed growth and compare
the resulting hierarchy with higher-degree monodromy for higher-order
Fueter maps. One may also study bounded inverse operators on weighted
Sobolev spaces after projecting out the periods. The current compact-set
path estimates do not claim such a global Sobolev theory. These are
questions left open by this draft, not prerequisites for its results.

\appendix
\section{The two distributional facts used in the proofs}
\label{app:distribution}

\subsection{Distributions supported at one point}
\label{app:point}
\begin{lemma}\label{lem:point-support}
A distribution $T$ supported at $0\in\R^n$ is a finite linear
combination of derivatives of $\delta_0$.
\end{lemma}
\begin{proof}
Choose a compact neighborhood $K$ of zero and an order $r$ and constant
$C$ such that
\[
 |\langle T,\varphi\rangle|
 \le C\max_{|\alpha|\le r}\sup_K|\partial^\alpha\varphi|
\]
for tests supported in $K$. Such a finite-order estimate on a fixed
compact set is part of the continuity definition of a distribution.
Let $\chi$ be a smooth cutoff equal to one near zero and set
$\chi_\varepsilon(x)=\chi(x/\varepsilon)$. Because the support of $T$
is zero, $\langle T,\varphi\rangle=
\langle T,\chi_\varepsilon\varphi\rangle$.

Suppose all derivatives of $\varphi$ through order $r$ vanish at zero.
Taylor's theorem for each derivative gives
\[
 \sup_{|x|\le C_0\varepsilon}|\partial^\beta\varphi(x)|
       =o(\varepsilon^{r-|\beta|}),\qquad |\beta|\le r.
\]
The product rule and the bounds for derivatives of
$\chi_\varepsilon$ then give
$\max_{|\alpha|\le r}\sup|
\partial^\alpha(\chi_\varepsilon\varphi)|=o(1)$.
The finite-order estimate forces $\langle T,\varphi\rangle=0$.
Therefore $T$ depends only on the finite $r$-jet of its test function
at zero. Linear functionals on that finite-dimensional jet space are
linear combinations of derivative evaluations, which are precisely
pairings with derivatives of $\delta_0$, up to their usual signs.
The vector-valued statement follows componentwise.
\end{proof}

The scaling argument in Lemma~\ref{lem:defect} uses more than this
structure theorem: local $L^1$ integrability and the first-order
operator eliminate all positive-order derivatives. This is the step
that makes the residue classification exhaustive rather than a list
of examples.

\subsection{An \texorpdfstring{$L^1$}{L1} Weyl lemma}
\label{app:Weyl}
\begin{lemma}\label{lem:Weyl}
If $u\in L^1_{\mathrm{loc}}(V)$ on an open $V\subset\R^n$ and
$\Delta u=0$ distributionally, then $u$ has a smooth harmonic
representative.
\end{lemma}
\begin{proof}
Convolve with a smooth compactly supported approximate identity to
obtain $u_\varepsilon$ on smaller subdomains. It is smooth and harmonic
because the Laplacian commutes with convolution. Its local $L^1$ norms
are uniformly bounded by the $L^1$ norm of $u$ on a slightly larger
compact set.

For a smooth harmonic function $h$, average its ordinary spherical
mean-value identities against a fixed smooth radial weight of total
integral one, supported in a ball of radius $r$. This writes $h$ in
a smaller ball as convolution with a smooth kernel $\psi_r$.
Differentiating the kernel gives, for every multi-index $\alpha$,
\[
 \sup_{B'}|\partial^\alpha h|
 \le C_{\alpha,r}\int_B|h(x)|\dd x,
\]
where $B'$ is a concentric smaller ball and the kernel supports fit
inside $B$. Apply this to $u_\varepsilon$. The uniform $L^1$ bounds
produce uniform bounds on all derivatives on compact subsets.
Repeated Arzel\`a--Ascoli extraction gives a subsequence converging
smoothly on compact subsets to a harmonic function $h$.
Since $u_\varepsilon\to u$ in $L^1_{\mathrm{loc}}$, the limit is a
representative of $u$. Local representatives agree almost everywhere
and hence everywhere by continuity, so they give the desired global
representative.
\end{proof}

For quaternion-valued fields the lemma applies to each of four real
components. The factorization $\overline\D\D=\lap$ then upgrades an
$L^1$ distributionally monogenic field to a smooth monogenic field,
which is the only elliptic regularity input required above.

\section{Reproduction and normalization checklist}
\label{app:reproduce}
The article uses the positive Laplacian
$\lap=\sum_{\mu=0}^3\partial_{x_\mu}^2$ and the left Fueter operator
$\D$. A convention using $-\lap$ changes the signs of the model fields
and their residues. It does not change the squared energy coefficient.
Our fundamental kernel contains $1/(2\pi^2)$, the sphere measure on
$\SSS$ has total area $4\pi$, and all planar loops are
counterclockwise-positive. These four choices account for the constants
$2/v$, $8\pi v$, and $8$ in the source, flux, and energy formulas.

The archive contains the self-contained LaTeX source, its compiled PDF,
\texttt{verify.py}, \texttt{verification\_results.json}, a captured
verification log, \texttt{requirements.txt}, and a README. No external
figure files or bibliography database are needed. Compile the source
with \texttt{pdflatex} twice to resolve references. Run the verification
script in a Python environment with SymPy, NumPy, and mpmath installed.
The exact package versions used for the reported run are recorded in
the JSON output and README.

\begin{thebibliography}{99}
\addcontentsline{toc}{section}{References}

\bibitem{GP}
R.~Ghiloni and A.~Perotti,
\emph{Slice regular functions on real alternative algebras},
Advances in Mathematics \textbf{226} (2011), 1662--1691.
\href{https://arxiv.org/abs/1008.4318}{arXiv:1008.4318}.

\bibitem{CSS}
F.~Colombo, I.~Sabadini, and F.~Sommen,
\emph{The inverse Fueter mapping theorem},
Communications on Pure and Applied Analysis \textbf{10} (2011), 1165--1181.
\href{https://doi.org/10.3934/cpaa.2011.10.1165}{doi:10.3934/cpaa.2011.10.1165}.

\bibitem{CPSS}
F.~Colombo, D.~Pe\~na Pe\~na, I.~Sabadini, and F.~Sommen,
\emph{A new integral formula for the inverse Fueter mapping theorem},
Journal of Mathematical Analysis and Applications \textbf{417} (2014), 112--122.
\href{https://arxiv.org/abs/1302.0685}{arXiv:1302.0685}.
The preprint's Theorem~3 is the rectangular-domain inverse result
explicitly inspected for this audit.

\bibitem{DQ}
B.~Dong and T.~Qian,
\emph{Further aspects of the Fueter mapping theorem},
preprint (2018).
\href{https://arxiv.org/abs/1805.02966}{arXiv:1805.02966}.
Theorem~3.6 gives a contour formulation of surjectivity.

\bibitem{DMP}
A.~De Martino and S.~Pinton,
\emph{A variation of the inverse Fueter theorem and the generalized
polyanalytic Cauchy--Kovalevskaya extension of order 2},
preprint, 7 August 2026.
\href{https://arxiv.org/abs/2608.07711}{arXiv:2608.07711}.

\bibitem{Common}
A.~K.~Common and F.~Sommen,
\emph{Axial monogenic functions from holomorphic functions},
Journal of Mathematical Analysis and Applications \textbf{179} (1993), 610--629.
\href{https://doi.org/10.1006/jmaa.1993.1372}{doi:10.1006/jmaa.1993.1372}.
Bibliographic lead; complete text not inspected in this audit.

\bibitem{DKQS}
B.~Dong, K.~I.~Kou, T.~Qian, and I.~Sabadini,
\emph{On the inversion of Fueter's theorem},
Journal of Geometry and Physics \textbf{108} (2016), 102--116.
\href{https://www.sciencedirect.com/science/article/pii/S0393044016301401}{Publisher record}.
Bibliographic lead; complete text not inspected in this audit.

\bibitem{Uploaded}
\emph{Classical Complex Theorems over the Quaternions:
Slice-regular and Cauchy--Fueter analogues, with proofs},
unpublished expository manuscript supplied with the request,
20 September 2026, file \texttt{quaternionic\_analysis.tex}.
See the section ``Fueter's mapping theorem and a constructive local
inverse,'' particularly the constructive axial inverse and its period caveat.

\end{thebibliography}
\end{document}
